\documentclass[10pt]{article}
\usepackage[utf8]{inputenc}
\usepackage[T1]{fontenc}
\usepackage{amsmath}
\usepackage{amsfonts}
\usepackage{amssymb}
\usepackage[version=4]{mhchem}
\usepackage{stmaryrd}
\usepackage{hyperref}
\hypersetup{colorlinks=true, linkcolor=blue, filecolor=magenta, urlcolor=cyan,}
\urlstyle{same}
\usepackage{bbold}

\title{A few more quadratic APN functions }

\author{Carl Bracken • Eimear Byrne • Nadya Markin • Gary McGuire}
\date{}


%New command to display footnote whose markers will always be hidden
\let\svthefootnote\thefootnote
\newcommand\blfootnotetext[1]{%
  \let\thefootnote\relax\footnote{#1}%
  \addtocounter{footnote}{-1}%
  \let\thefootnote\svthefootnote%
}

%Overriding the \footnotetext command to hide the marker if its value is `0`
\let\svfootnotetext\footnotetext
\renewcommand\footnotetext[2][?]{%
  \if\relax#1\relax%
    \ifnum\value{footnote}=0\blfootnotetext{#2}\else\svfootnotetext{#2}\fi%
  \else%
    \if?#1\ifnum\value{footnote}=0\blfootnotetext{#2}\else\svfootnotetext{#2}\fi%
    \else\svfootnotetext[#1]{#2}\fi%
  \fi
}

\begin{document}
\maketitle
Received: 26 April 2008 / Accepted: 27 October 2010 / Published online: 11 November 2010\\
© Springer Science+Business Media, LLC 2010

\begin{abstract}
APN functions on $G F\left(2^{n}\right)$ where $n$ is divisible by 3 but not 9 . The family contains inequivalent functions, obtained by setting some coefficients equal to 0 . We also discuss the inequivalence proof (by computation) which shows that these functions are new.
\end{abstract}

Keywords Almost perfect nonlinear • Finite field • APN • CCZ equivalent • EA equivalent

Mathematics Subject Classifications (2010) 11T06•11T71

\section*{1 Introduction}
Let $L=G F\left(2^{n}\right)$ for some positive integer $n$. A function $f: L \longrightarrow L$ is said to be almost perfect nonlinear (APN) on $L$ if the number of solutions in $L$ of the equation

$$
f(x+q)+f(x)=p
$$

is at most 2 , for all $p, q \in L, q \neq 0$. Equivalently, $f$ is APN if the set $\{f(x+q)+ f(x): x \in L\}$ has size at least $2^{n-1}$ for each $q \in L^{*}$. Clearly, as $L$ has characteristic 2 , the number of solutions to the above equation must be an even number for any function $f$ on $L$.

APN functions were introduced in [13] by Nyberg, who defined them as the mappings with highest resistance to differential cryptanalysis. In other words, APN

\footnotetext{Research of Carl Bracken was supported by Irish Research Council for Science, Engineering and Technology Postdoctoral Fellowship.

Research of Eimear Byrne, Nadya Markin and Gary McGuire was supported by the Claude Shannon Institute, Science Foundation Ireland Grant 06/MI/006.\\
C. Bracken . E. Byrne . N. Markin . G. McGuire ( $\boxtimes$ )

School of Mathematical Sciences, University College Dublin, Dublin, Ireland\\
e-mail: \href{mailto:gary.mcguire@ucd.ie}{gary.mcguire@ucd.ie}
}
functions are those for which the plaintext difference $x-y$ yields the ciphertext difference $f(x)-f(y)$ with probability $1 / 2^{n-1}$. Since Nyberg's characterization, many papers have been written on APN functions, although not many different families of such functions are known.

The main result of this paper is a construction of a new family of APN functions, in Theorem 2.1.

Two functions $f, g: L \longrightarrow L$ are called extended affine (EA) equivalent if there exist affine permutations $A_{1}, A_{2}$ and an affine map $A$ such that $g=A_{1} \circ f \circ A_{2}+A$.

Until recently, all known APN functions were EA-equivalent to one of a short list of monomial functions, namely the Gold functions $\left(f(x)=x^{2^{k}+1}\right.$ where $\left.(k, n)=1\right)$, Kasami-Welch functions $\left(f(x)=x^{4^{k}-2^{k}+1}\right.$ where $\left.(k, n)=1\right)$, inverse, Welch, Niho and Dobbertin functions. The last four are APN only when $n$ is odd, we do not need them for this paper. For some time it was conjectured that this list was the complete list of APN functions up to EA equivalence. It is still conjectured to be the complete list of monomial APN functions, up to CCZ equivalence.

This more general notion of equivalence has been suggested in [9], and is referred to as Carlet-Charpin-Zinoviev (CCZ) equivalence. Two functions are called CCZequivalent if the graph of one can be obtained from the graph of the other by an affine permutation of the product space. EA equivalence is a special case of CCZ equivalence.

We say that $f: L \longrightarrow L$ is differentially $m$-uniform if the polynomial $f(x+q)+ f(x)+p$ has at most $m$ roots in $L$, for any $p, q \in L, q \neq 0$. Then $f$ is APN on $L$ if and only if it is differentially 2-uniform on $L$. Differential uniformity, and resistance to linear and differential attacks, are invariants of CCZ-equivalence.

In [5], Proposition 3, the authors express necessary and sufficient conditions for EA equivalence of functions in terms of CCZ-equivalence and use this to construct several examples of APN functions that are CCZ-equivalent to the Gold functions, but not EA-equivalent to any monomial function. This showed that the original conjecture is false. The new question was whether all APN functions are CCZequivalent to one on the list.

In 2006 a sporadic example of a binomial APN function that is not CCZ equivalent to any power mapping was given in [11]. A family of APN binomials on fields $\mathbb{F}_{2^{n}}$, where $n$ is divisible by 3 but not 9 , was presented in [4]. In [7] these have been shown to be EA inequivalent to any monomial function, and CCZ inequivalent to the Gold or Kasami-Welch functions. For the case $n=6$, in [10] Dillon presented a list of CCZ inequivalent APN functions on $G F\left(2^{n}\right)$, found by computer search.

Below we list all the infinite families of non-monomial APN functions known at the time of writing. These families are all pairwise CCZ inequivalent.\\
1.

$$
f(x)=x^{2^{s}+1}+\alpha x^{2^{i k}+2^{m k+s}}
$$

where $n=3 k,(k, 3)=(s, 3 k)=1, k \geq 3, i \equiv s k \bmod 3, m \equiv-i \bmod 3, \alpha= t^{2^{k}-1}$ and $t$ is primitive (see Budaghyan et al. [4]).\\
2.

$$
f(x)=x^{2^{s}+1}+\alpha x^{2^{i k}+2^{m k+s}}
$$

where $n=4 k,(k, 2)=(s, 2 k)=1, k \geq 3, i \equiv s k \bmod 4, m=4-i, \alpha=t^{2^{k}-1}$ and $t$ is primitive (see Budaghyan et al. [6]). This family generalizes an example found for $n=12$ by Edel et al. [11].\\
3.

$$
f(x)=\alpha x^{2^{s}+1}+\alpha^{2^{k}} x^{2^{k+s}+2^{k}}+\beta x^{2^{k}+1}+\sum_{i=1}^{k-1} \gamma_{i} x^{2^{k+i}+2^{i}},
$$

where $n=2 k, \alpha$ is not a cube, $\beta \notin G F\left(2^{k}\right), \gamma_{i} \in G F\left(2^{k}\right)$ for each $i$, and $(k, s)=1$, $k$ is odd, $s$ is odd (see Bracken et al. [1]).\\
4.

$$
F(x)=x^{2^{2 s}+2^{s}}+b x^{2^{k}+1}+c x^{2^{k}\left(2^{2 s}+2^{s}\right)}
$$

where $n=2 k,(s, k)=1,\left(2^{s}+1,2^{k}+1\right)>1, c^{2^{k}+1}=1, c b^{2^{k}} \neq b$ and $c \notin \left\{\lambda^{\left(2^{s}+1\right)\left(2^{k}-1\right)}: \lambda \in G F\left(2^{2 k}\right)\right\}$ (see Budaghyan and Carlet [3]).\\
5.

$$
F(x)=\theta x^{2^{s+k}+1}+\theta^{2^{k}} x^{2^{k}+2^{s}}+x^{2^{s}+1}+x^{2^{k}\left(2^{s}+1\right)}+x^{2^{k}+1}+\pi x^{2^{s+k}+2^{s}}
$$

where $n=2 k,(s, k)=1, \pi \notin G F\left(2^{k}\right)$ and the equation $x^{2^{s}+1}+\theta x^{2^{s}}+\theta^{2^{k}} x+1=0$ has no solution $x$ satisfying $x^{2^{k}+1}=1$ [3].\\
6.

$$
f(x)=x^{3}+\operatorname{Tr}\left(x^{9}\right)
$$

over $G F\left(2^{n}\right)$, any $n$ (see Budaghyan et al. [8]).\\
7.

$$
f(x)=u x^{2^{-k}+2^{k+s}}+u^{2^{k}} x^{2^{s}+1}+v x^{2^{k+s}+2^{s}}
$$

where $n=3 k, u$ is primitive, $v \in G F\left(2^{k}\right),(s, 3 k)=1,(3, k)=1$ and 3 divides $k+s$ (see Bracken et al. [1]).\\
8.

$$
F(x)=u^{2^{k}} x^{2^{-k}+2^{k+s}}+u x^{2^{s}+1}+v x^{2^{-k}+1}+w u^{2^{k}+1} x^{2^{k+s}+2^{s}}
$$

where $n=3 k, s$ and $k$ are positive integers with $k+s$ divisible by three and $(s, 3 k)=(3, k)=1, u$ is a primitive element of $G F\left(2^{3 k}\right)$ and $v, w \in G F\left(2^{k}\right)$ with $v w \neq 1$ (this paper).

In general, establishing CCZ equivalence of arbitrary functions is extremely difficult. There are, however, a number of invariants of CCZ equivalence that can be useful in the classification of functions. A nice link with coding theory is that a pair of functions $f$ and $g$ on $L$ are CCZ equivalent on $L$ if and only if the binary codes with parity check matrices

$$
H_{f}=\left[\begin{array}{ccc}
1 & \cdots & 1 \\
\mathbf{x}_{1} & \cdots & \mathbf{x}_{2^{n}} \\
f\left(\mathbf{x}_{1}\right) & \cdots & f\left(\mathbf{x}_{2^{n}}\right)
\end{array}\right], \quad H_{g}=\left[\begin{array}{ccc}
1 & \cdots & 1 \\
\mathbf{x}_{1} & \cdots & \mathbf{x}_{2^{n}} \\
g(\mathbf{x}) & \cdots & g\left(\mathbf{x}_{2^{n}}\right)
\end{array}\right]
$$

are equivalent over $G F(2)$ (that is one can be obtained from the other by a monomial transformation), see [1]. Here $\mathbf{x}_{i}, f\left(\mathbf{x}_{i}\right)$ and $g\left(\mathbf{x}_{i}\right)$ are expressions of $x_{i}, f\left(x_{i}\right)$ and $g\left(x_{i}\right)$\\
respectively as binary vectors of length $n$ in $L$ viewed as a $G F(2)$ vector space and $L=\left\{x_{1}, \ldots, x_{2^{n}}\right\}$.

In this paper we introduce a new family of APN functions on fields of order $2^{3 k}$ where $k$ is not divisible by 3 . The family of polynomials has the form

$$
F(x)=u^{2^{k}} x^{2^{-k}+2^{k+s}}+u x^{2^{s}+1}+v x^{2^{-k}+1}+w u^{2^{k}+1} x^{2^{k+s}+2^{s}}
$$

with certain constraints on the integers $s, k$ and on $u, v, w \in G F\left(2^{3 k}\right)$ (see Theorem 2.1).

Curiously, setting $v=0$ gives a family of trinomial APN functions (Family 7) that contains inequivalent functions to some of the functions with $v \neq 0$. Setting $w=0$ gives yet another family with inequivalent functions. We discuss this in Section 3.

The layout of this paper is as follows. In the next section we show that our polynomials are indeed APN functions on $G F\left(2^{3 k}\right)$. Using code equivalence, in Section 3 we discuss the fact that these functions are not CCZ equivalent to any known APN functions when $n=12$, and are therefore new.

\section*{2 New APN functions}
The following theorem gives a construction of quadratic quadrinomial APN functions on $G F\left(2^{n}\right)$ whenever $n$ is divisible by 3 but not 9 . A quadratic monomial is one of the form $x^{2^{i}+2^{j}}$ for some integers $i$ and $j$. Observe that if $f(x)=x^{2^{i}+2^{j}}$, then

$$
f(x+q)+f(x)+f(q)=x^{2^{i}} q^{2^{j}}+x^{2^{j}} q^{2^{i}}
$$

is a linear function in $x$, whose kernel has the same size as the set of zeroes of any of its translates, such as the solution set of $f(x)+f(x+q)=p$ in $L$, for any $p \in L$. Because of this property, proving whether or not a quadratic polynomial is APN is more tangible than one that is not quadratic. For this reason, all of the recently discovered families of APN functions have been quadratic.

We will show that our polynomial $F(x)$ is APN by computing the size of the kernel of the corresponding linear map

$$
F(x+q)+F(x)+F(q)
$$

Theorem 2.1 Let $s$ and $k$ be positive integers with $k+s$ divisible by three and $(s, 3 k)= (3, k)=1$. Let $u$ be a primitive element of $G F\left(2^{3 k}\right)$ and let $v, w \in G F\left(2^{k}\right)$ with $v w \neq 1$. Then the function

$$
F(x)=u^{2^{k}} x^{2^{-k}+2^{k+s}}+u x^{2^{s}+1}+v x^{2^{-k}+1}+w u^{2^{k}+1} x^{2^{k+s}+2^{s}}
$$

is APN over $G F\left(2^{3 k}\right)$.\\
Proof We show that for every $p$ and $q$ (with $q \neq 0$ ) in $G F\left(2^{3 k}\right)$ the equation

$$
F(x)+F(x+q)=p
$$

has at most two solutions by counting the number of solutions to the equation

$$
F(x)+F(x+q)+F(q)=0
$$

This gives

$$
\begin{aligned}
F(x)+F(x+q)+F(q)= & u^{2^{k}}\left(x^{2^{k+s}} q^{2^{-k}}+q^{2^{k+s}} x^{2^{-k}}\right)+u\left(x^{2^{s}} q+q^{2^{s}} x\right) \\
& +v\left(x^{2^{-k}} q+q^{2^{-k}} x\right)+w u^{2^{k}+1}\left(x^{2^{s}} q^{2^{k+s}}+q^{2^{s}} x^{2^{k+s}}\right)=0
\end{aligned}
$$

Replace $x$ with $x q$ to obtain

$$
\begin{aligned}
u^{2^{k}} q^{2^{-k}+2^{k+s}}\left(x^{2^{k+s}}+x^{2^{-k}}\right) & +u q^{2^{s}+1}\left(x^{2^{s}}+x\right)+v q^{2^{-k}+1}\left(x^{2^{-k}}+x\right) \\
& +w u^{2^{k}+1} q^{2^{k+s}+2^{s}}\left(x^{2^{s}}+x^{2^{k+s}}\right)=0
\end{aligned}
$$

and collect terms in $x$ to get

$$
\begin{aligned}
\Delta(x):= & \left(v q^{2^{-k}+1}+u q^{2^{s}+1}\right) x+\left(v q^{2^{-k}+1}+u^{2^{k}} q^{2^{-k}+2^{k+s}}\right) x^{2^{-k}} \\
& +\left(w u^{2^{k}+1} q^{2^{k+s}+2^{s}}+u q^{2^{s}+1}\right) x^{2^{s}}+\left(w u^{2^{k}+1} q^{2^{k+s}+2^{s}}\right. \\
& \left.+u^{2^{k}} q^{2^{-k}+2^{k+s}}\right) x^{2^{k+s}}=0
\end{aligned}
$$

We write

$$
\Delta(x)=A x+B x^{2^{-k}}+C x^{2^{s}}+D x^{2^{k+s}}
$$

where

$$
\begin{aligned}
& A=v q^{2^{-k}+1}+u q^{2^{s}+1}, \quad B=v q^{2^{-k}+1}+u^{2^{k}} q^{2^{-k}+2^{k+s}} \\
& C=w u^{2^{k}+1} q^{2^{k+s}+2^{s}}+u q^{2^{s}+1}, \quad D=w u^{2^{k}+1} q^{2^{k+s}+2^{s}}+u^{2^{k}} q^{2^{-k}+2^{k+s}}
\end{aligned}
$$

Clearly 0 is a root of $\Delta(x)$. Moreover $\Delta(1)=A+B+C+D=0$. If we show that 0 and 1 are the only solutions of $\Delta(x)=0$, then we will have proved that $F(x)$ is APN on $G F\left(2^{3 k}\right)$.

First we demonstrate that none of $A, B, C$ or $D$ vanish for any $q \in G F\left(2^{3 k}\right)^{*}$. If $A=0$ we have $u=v q^{2^{-k}-2^{s}}$ which implies $u^{2^{k}}=v q^{1-2^{k+s}}$. By hypothesis, $k+s$ is divisible by 3 , so that $1-2^{k+s}$ is divisible by 7 , and hence $q^{1-2^{k+s}}$ is a 7 th power in $G F\left(2^{3 k}\right)$. Since 3 does not divide $k, 7$ does not divide $2^{k}-1$, so the map $x \mapsto x^{7}$ is a permutation on $G F\left(2^{k}\right)$. Then $v \in G F\left(2^{k}\right)$ can be expressed as a 7th power. This means that $u^{2^{k}}$ and hence $u$ is a 7th power in $G F\left(2^{3 k}\right)$. This gives a contradiction because 7 is a divisor of $2^{3 k}-1$ and we chose $u$ to be primitive in $G F\left(2^{3 k}\right)$. We deduce that $A \neq 0$. Similar arguments show that $B, C$ and $D$ are all nonzero.

Next we define the linearized polynomial:

$$
L_{\theta}(T):=T+\theta T^{2^{k}}+\theta^{2^{k}+1} T^{2^{-k}}
$$

When $T=\theta x+x^{2^{-k}}$ and $\theta$ is a ( $2^{k}-1$ )-th power, a routine calculation verifies that $L_{\theta}(T)=0$ for all $x \in G F\left(2^{3 k}\right)$. Observe that

$$
\frac{A}{B}=\frac{v q^{2^{-k}+1}+u q^{2^{s}+1}}{v q^{2^{-k}+1}+u^{2^{k}} q^{2^{-k}+2^{k+s}}}=\frac{v+u q^{2^{s}-2^{-k}}}{v+u^{2^{k}} q^{2^{k+s}-1}}=\left(v+u q^{2^{s}-2^{-k}}\right)^{1-2^{k}}
$$

which gives


\begin{equation*}
L_{\frac{A}{B}}\left(\frac{A}{B} x+x^{2^{-k}}\right)=0 \tag{1}
\end{equation*}


Suppose $\Delta(x)=0$ and consider the equation

$$
\frac{\Delta(x)}{B}=\left(\frac{A}{B} x+x^{2^{-k}}\right)+\left(\frac{C}{B} x^{2^{s}}+\frac{D}{B} x^{2^{k+s}}\right)=0 .
$$

Applying this to (1) gives

$$
L_{\frac{A}{B}}\left(\frac{\Delta(x)}{B}\right)=L_{\frac{A}{B}}\left(\frac{C}{B} x^{2^{s}}+\frac{D}{B} x^{2^{k+s}}\right)=0 .
$$

We compute this as

$$
\begin{aligned}
\left(B^{2^{-k}+2^{k}} C\right. & \left.+D^{2^{-k}} A^{2^{k}+1}\right) x^{2^{s}}+\left(B^{2^{-k}+2^{k}} D+B^{2^{-k}} A C^{2^{k}}\right) x^{2^{k+s}} \\
& +\left(B^{2^{-k}} A D^{2^{k}}+A^{2^{k}+1} C^{2^{-k}}\right) x^{2^{-k+s}}=0
\end{aligned}
$$

We substitute in the values of $A, B, C$, and $D$ and after simplification we obtain the following

$$
\begin{aligned}
& (v w+1) u q^{2^{k}}+1+2^{s}\left(v q^{2^{-k}}+u q^{2^{s}}\right)\left(u^{2^{k}} q^{2^{k+s}+2^{k}}+u^{2^{-k}} q^{2^{-k+s}+1}\right) x^{2^{s}} \\
& \quad+(v w+1) u^{2^{k}} q^{2^{k}+1+2^{k+s}}\left(v q^{2^{-k}}+u q^{2^{s}}\right)\left(u q^{2^{k}+2^{s}}+u^{2^{-k}} q^{2^{-k+s}+2^{-k}}\right) x^{2^{k+s}} \\
& \quad+(v w+1) u^{2^{-k}} q^{2^{k}+1+2^{-k+s}}\left(v q^{2^{-k}}+u q^{2^{s}}\right)\left(u^{2^{k}} q^{2^{k+s}+2^{-k}}+u q^{2^{s}+1}\right) x^{2^{-k+s}}=0
\end{aligned}
$$

As we chose $v$ and $w$ such that $v \neq w^{-1}$ and as $A \neq 0$ we can divide the equation by $(v w+1) q^{2^{k}+1}\left(v q^{2^{-k}}+u q^{2^{s}}\right) u^{2^{-k}+1} q^{2^{-k+s}+2^{s}+1}$ and take the expression to the $2^{-s}$ power to obtain


\begin{equation*}
\left(1+a^{-2^{k-s}}\right) x+\left(a^{2^{-s}}+a^{-2^{k-s}}\right) x^{2^{k}}+\left(1+a^{2^{-s}}\right) x^{2^{-k}}=0 \tag{2}
\end{equation*}


where $a=u^{2^{k}-1} q^{2^{-k}+2^{k+s}-2^{s}-1}$.\\
Recalling that $\Delta(x)=0$ by assumption, we have $L_{\frac{D}{C}}\left(\frac{\Delta(x)}{D}\right)=0$. Moreover, we know $L_{\frac{D}{C}}\left(x^{2^{s}}+\frac{D}{C} x^{2^{k+s}}\right)=0$, as

$$
\frac{D}{C}=\frac{w u^{2^{k}+1} q^{2^{k+s}+2^{s}}+u^{2^{k}} q^{2^{-k}+2^{k+s}}}{w u^{2^{k}+1} q^{2^{k+s}+2^{s}}+u q^{2^{s}+1}}=\left(w+u^{-1} q^{2^{-k}-2^{s}}\right)^{1-2^{k}}
$$

This implies $L_{\frac{D}{C}}\left(\frac{A}{C} x+\frac{B}{C} x^{2^{-k}}\right)=0$, which we compute as

$$
\begin{aligned}
\left(C^{2^{-k}+2^{k}} A\right. & \left.+C^{2^{-k}} D B^{2^{k}}\right) x+\left(C^{2^{-k}} D A^{2^{k}}+D^{2^{k}+1} B^{2^{-k}}\right) x^{2^{k}} \\
& +\left(C^{2^{-k}+2^{k}} B+D^{2^{k}+1} A^{2^{-k}}\right) x^{2^{-k}}=0
\end{aligned}
$$

This yields

$$
\begin{aligned}
& u^{2^{k}+2^{-k}} q^{2^{s+k}+2^{s-k}+1}\left(w u q^{2^{s}}+q^{2^{-k}}\right)(1+w v)\left(u q^{2^{s}+2^{k}}+u^{2^{-k}} q^{2^{-k+s}+2^{-k}}\right) x \\
& \quad+u^{2^{k}+2^{-k}} q^{2^{s+k}+2^{s-k}+2^{k}}\left(w u q^{2^{s}}+q^{2^{-k}}\right)(1+w v)\left(u q^{2^{s}+1}+u^{2^{k}} q^{2^{-k}+2^{k+s}}\right) x^{2^{k}} \\
& \quad+u^{2^{k}+2^{-k}} q^{2^{s+k}+2^{s-k}+2^{-k}}\left(w u q^{2^{s}}+q^{2^{-k}}\right)(1+w v)\left(u^{2^{-k}} q^{1+2^{-k+s}}+u^{2^{k}} q^{2^{k+s}+2^{k}}\right) x^{2^{-k}}
\end{aligned}
$$

Noting that $C=u q^{2^{s}}\left(w u q^{2^{s}}+q^{2^{-k}}\right)^{2^{k}} \neq 0$, we may divide the above equation by $u^{2^{k}+2^{-k}+1} q^{2^{s+k}+2^{s-k}+2^{k}+2^{s}+1}\left(w u q^{2^{s}}+q^{2^{-k}}\right)(1+w v)$ to obtain:


\begin{equation*}
\left(1+a^{-2^{-k}}\right) x+(1+a) x^{2^{k}}+\left(a+a^{-2^{-k}}\right) x^{2^{-k}}=0 \tag{3}
\end{equation*}


Now we combine (2) and (3) such that the terms in $x^{2^{-k}}$ cancel. This will give

$$
\begin{aligned}
& \left(\left(1+a^{-2^{k-s}}\right)\left(a+a^{-2^{-k}}\right)+\left(1+a^{-2^{-k}}\right)\left(1+a^{2^{-s}}\right)\right) x \\
& \quad+\left(\left(a^{2^{-s}}+a^{-2^{k-s}}\right)\left(a+a^{-2^{-k}}\right)+(1+a)\left(1+a^{2^{-s}}\right)\right) x^{2^{k}}=0
\end{aligned}
$$

which is the same as

$$
\left(\left(1+a^{-2^{k-s}}\right)\left(a+a^{-2^{-k}}\right)+\left(1+a^{-2^{-k}}\right)\left(1+a^{2^{-s}}\right)\right)\left(x+x^{2^{k}}\right)=0 .
$$

If we show that $\left(1+a^{2^{k-s}}\right)\left(a+a^{-2^{-k}}\right)+\left(1+a^{-2^{-k}}\right)\left(1+a^{2^{-s}}\right) \neq 0$ for all possible values of $a$ then we could conclude that $x \in G F\left(2^{k}\right)$. To this end we consider the expression

$$
\left(1+a^{-2^{k-s}}\right)\left(a+a^{-2^{-k}}\right)=\left(1+a^{-2^{-k}}\right)\left(1+a^{2^{-s}}\right)
$$

Rearranging we obtain, since $a$ is a ( $2^{k}-1$ )-th power,

$$
a=\frac{\left(1+a^{-1}\right)^{2^{-k}}}{\left(1+a^{-1}\right)^{2^{k-s}}} \frac{(1+a)^{2^{-s}}}{(1+a)^{2^{k}}} .
$$

This implies $a$ is a $\left(2^{k+s}-1\right)$-th power because

$$
a=\frac{\left(1+a^{-1}\right)^{2^{-k}}}{\left(1+a^{-1}\right)^{2^{k-s}}} \frac{(1+a)^{2^{-s}}}{(1+a)^{2^{k}}}=\left(1+a^{-1}\right)^{2^{k-s}\left(2^{k+s}-1\right)}(1+a)^{2^{-s}\left(1-2^{k+s}\right)} .
$$

This in turn implies that $a$ is a seventh power. As $a=u^{2^{k}-1} q^{2^{-k}+2^{k+s}-2^{s}-1}= u^{2^{k}-1} q^{\left(2^{k+s}-1\right)\left(1-2^{-k}\right)}$ we see that if $a$ is a seventh power then so is $u^{2^{k}-1}$ but this is not possible because $k$ is not divisible by three and $u$ is primitive. We can now state that all solutions to $\Delta(x)=0$ are in $G F\left(2^{k}\right)$. Applying this to our original expression for $\Delta(x)$ gives

$$
\left(u q^{2^{s}+1}+u^{2^{k}} q^{2^{-k}+2^{k+s}}\right)\left(x+x^{2^{s}}\right)=0
$$

If $u q^{2^{s}+1}+u^{2^{k}} q^{2^{-k}+2^{k+s}}=0$ then $a=1$ because $a=u^{2^{k}-1} q^{2^{-k}+2^{k+s}-2^{s}-1}$. But 1 is a seventh power, so we must have $a \neq 1$. Then $x+x^{2^{s}}=0$, which implies $x=0$ or 1 because $s$ is relatively prime to $3 k$.

\section*{3 Equivalence}
It remains to show that the "new" family of APN functions introduced in this paper is indeed new. We therefore need to demonstrate that these functions are not CCZ equivalent to any known APN function. The only known method of proving CCZ inequivalence without a computer involves a brute force approach, and is only feasible for power functions. We resort here to a demonstration by computer for small values of $n$, which is commonly done, see [3] for example. We attempt to show that the corresponding error-correcting codes are inequivalent, which is necessary and sufficient as we said in the introduction, and is proved in [1].

First we wish to point out that our Family 8 has functions when the degree of the field $n$ is even, and hence cannot be equivalent to a function that is only defined on odd degree fields. This implies that our Family 8 is not CCZ equivalent to any of the monomial functions inverse, Welch, Niho, Dobbertin, because these are only APN on odd degree fields. Further, it should be noted that the power mappings are defined on fields not divisible by 3 , so they can describe functions not covered by our function.

A similar comment applies to Families 3 and 4. Since Families 3 and 4 are defined for $n$ divisible by 2 but not by 4 , they cannot completely cover Family 8 (which can be defined for $n$ divisible by 4 ).

One standard method of proving two codes to be inequivalent is to show that they have a different weight distribution (when this is the case). However, all the computer evidence shows that the codes from this paper all have the same weight distribution as the code for the function $x^{3}$ (and we have proved this for Family 7 in [2]). We will therefore use other invariants.

Two other invariants are called the gamma-rank and delta-rank. We shall give the definitions here, but see [10] and [11] for further discussion. Let $f$ be an APN function on $G F\left(2^{n}\right)$. Let $G_{f}=\left\{(x, f(x)): x \in G F\left(2^{n}\right)\right\} \subseteq G F\left(2^{n}\right) \times G F\left(2^{n}\right)$ be the graph of $f$. Let $D(f)$ denote the development of $G_{f}$, which is an incidence structure whose blocks are the translates of $G_{f}$ in $G F\left(2^{n}\right) \times G F\left(2^{n}\right)$. The incidence matrix of $D(f)$ is a $2^{2 n} \times 2^{2 n}$ matrix, and its rank over $G F(2)$ is called the gamma-rank of $f$.

The delta-rank of $f$ is the rank over $G F(2)$ of the $(0,1)$ matrix defined as follows: the rows and columns are labelled by the points of $G F\left(2^{n}\right) \times G F\left(2^{n}\right)$, and there is a 1 in position $\left(\left(x_{1}, y_{1}\right),\left(x_{2}, y_{2}\right)\right)$ if and only if there are two solutions to the equation $f(x)+f\left(x+x_{1}+x_{2}\right)=y_{1}+y_{2}$.\\
(Based on numerical evidence, we conjecture that all Gold functions have the same delta-rank.)

The representation of our quadrinomial Family 8

$$
F(x)=u^{2^{k}} x^{2^{-k}+2^{k+s}}+u x^{2^{s}+1}+v x^{2^{-k}+1}+w u^{2^{k}+1} x^{2^{k+s}+2^{s}}
$$

can be used to describe two of the families listed in the introduction. These are Families 1 and 7 (see the table below). There is also another family, of trinomials, obtained by setting $w=0$, which also appears to be inequivalent. Let us call this family (Family 8 with $w=0$ ) Family $8 W$.

\begin{center}
\begin{tabular}{|l|l|}
\hline
Function & Class \\
\hline
$u^{2^{k}} x^{2^{-k}+2^{k+s}}+u x^{2^{s}+1}+v x^{2^{-k}+1}+w u^{2^{k}+1} x^{2^{k+s}+2^{s}}$ & Thm 2.1, NEW (Family 8) \\
\hline
$u^{2^{k}} x^{2^{-k}+2^{k+s}}+u x^{2^{s}+1}+v x^{2^{-k}+1}$ & Thm 2.1, $v \neq 0, w=0$, NEW (Fam 8W) \\
\hline
$u^{2^{k}} x^{2^{-k}+2^{k+s}}+u x^{2^{s}+1}+w u^{2^{k}+1} x^{2^{k+s}+2^{s}}$ & Thm 2.1 with $v=0, w \neq 0$ (Family 7) \\
\hline
$u^{2^{k}} x^{2^{-k}+2^{k+s}}+u x^{2^{s}+1}$ & Thm 2.1, $v=w=0$ (Family 1) \\
\hline
\end{tabular}
\end{center}

We claim that these four families are pairwise CCZ inequivalent, and are also new (in the sense that they are not CCZ equivalent to any of the other families). As usual, we have no proof of this for general $n$, but we now present computer evidence for this claim.

For smaller dimensions than 12, CCZ equivalence can be directly determined by testing equivalence of the asssociated codes with Magma. For the case $n=6$ the polynomials in the four families under consideration take the following forms:

$$
\begin{aligned}
& u x^{3}+v u^{5} x^{10}+v x^{17}+u^{4} x^{24} \\
& u x^{3}+v x^{17}+u^{4} x^{24} \\
& u x^{3}+v u^{5} x^{10}+u^{4} x^{24} \\
& u x^{3}+u^{4} x^{24},
\end{aligned}
$$

for some primitive element $u \in G F\left(2^{6}\right)$ and $v \in G F(4)$. In the first three cases, the polynomials are CCZ equivalent to

$$
x^{3}+x^{10}+u x^{24},
$$

which appears in Dillon's list [10], and in the last case (the binomial) the polynomial is CCZ equivalent to $x^{3}$. Therefore, $n=6$ is not a sufficiently large value of $n$ to distinguish our four families, but does distinguish Family 1 from the others.

The next smallest possible value of $n$ to consider is $n=12$, so $k=4$. Example functions (with $s=5$ ) from the four families under consideration are in the following table.

\begin{center}
\begin{tabular}{ll}
\hline
Function & Class \\
\hline
$u^{16} x^{768}+u x^{33}+x^{257}+u^{290} x^{544}$ & Theorem 2.1, NEW (Family 8) \\
$u^{16} x^{768}+u x^{33}+x^{257}$ & Theorem 2.1, $v \neq 0, w=0$, NEW (Fam 8W) \\
$u^{16} x^{768}+u x^{33}+u^{290} x^{544}$ & Theorem 2.1 with $v=0, w \neq 0$ (Family 7) \\
$u^{16} x^{768}+u x^{33}$ & Theorem 2.1, $v=w=0$ (Family 1) \\
\hline
\end{tabular}
\end{center}

Magma has a built in test for code equivalence. This test involves performing a backtrack search using the action of the automorphism group of the code on the words of minimum weight. However, for $n=12$ each of these codes has $1,397,760$ words of miniumum weight and this is beyond the capability of the Leon package PERM for code equivalence that is used in Magma and other systems. So Magma's current coding theory programs cannot determine whether the codes for Families $7,8,8 \mathrm{~W}$ are equivalent-the calculation is too big when $n=12$.

Using a newly developed method, John Cannon, Gabi Nebe and Allan Steel of the Magma group computed (Cannon et al., personal communication) the delta-rank of the four APN functions given above, along with the delta ranks of functions from all the other families. Their results are shown in the table below. The first three functions in the table above were found to have delta-rank 7,900 while the fourth has delta rank 7,816. with different delta ranks. Hence the fourth APN function (Family 1) is CCZ inequivalent to the first three (Families $7,8,8 \mathrm{~W}$ ) which confirms what we already deduced from the $n=6$ case. However, the delta-rank will still not distinguish Families $7,8,8 \mathrm{~W}$. All other delta-ranks (for functions from other families) were different to 7,900 and 7,816, as the reader can see from the table below. Therefore, these delta-ranks show that all other APN functions are not equivalent to Families $7,8,8 \mathrm{~W}$. The only monomial functions that are defined on $n=12$ are the Gold functions $x^{3}$ and $x^{33}$, and one Kasami-Welch function $x^{993}$.

\begin{center}
\begin{tabular}{|l|l|l|}
\hline
Family & Function & Delta-Rank \\
\hline
Gold & $x^{3}$ & 7,550 \\
\hline
Gold & $x^{33}$ & 7,550 \\
\hline
Kasami-Welch & $x^{993}$ & 62,550 \\
\hline
1 & $u^{16} x^{768}+u x^{33}$ & 7,816 \\
\hline
2 & $x^{3}+u^{7} x^{528}$ & 7,822 \\
\hline
3 & Does not exist on $n=12$ &  \\
\hline
4 & Does not exist on $n=12$ &  \\
\hline
5 & $x^{3}+x^{65}+u x^{129}+u^{64} x^{66}+u^{3} x^{130}+x^{192}$ & 7,550 \\
\hline
6 & $x^{3}+\operatorname{Tr}\left(x^{3}\right)$ & 7,846 \\
\hline
7 & $u^{16} x^{768}+u x^{33}+u^{290} x^{544}$ & 7,900 \\
\hline
8W & $u^{16} x^{768}+u x^{33}+x^{257}$ & 7,900 \\
\hline
8 & $u^{16} x^{768}+u x^{33}+x^{257}+u^{290} x^{544}$ & 7,900 \\
\hline
\end{tabular}
\end{center}

The problem still remains of checking Families $7,8,8 \mathrm{~W}$ for equivalence to each other. This was done by the Magma group (Cannon et al., personal communication), using their own method and a certain invariant. An independent verification of this result was provided by Thomas Feulner, who calculated a new invariant: a canonical form for the generator matrices of the codes. Full details of that algorithm, and the outputs, are available at [12]. All three sample functions above had different canonical forms. Feulner found that the example function from Family 8 has a different automorphism group to the functions from Families 7 and 8 W . However the functions from 7 and 8 W (both are trinomials) have the same automorphism group, so this invariant does not distinguish them.

Finally, we remark that there are no known cases of functions from the same family being inequivalent. In other words, different choices of coefficients within a family (apart from setting some coefficients to 0 ) always seem to produce equivalent functions.

Acknowledgements We thank John Cannon, Gabriele Nebe, and Allan Steel for their work on APN functions and Magma. We are very grateful to Thomas Feulner for his computations, and for putting the results online. We thank the anonymous referees whose comments led to some improvements to this paper.

\section*{References}
\begin{enumerate}
  \item Bracken, C., Byrne, E., Markin, N., McGuire, G.: New families of quadratic almost perfect nonlinear trinomials and multinomials. Finite Fields Their Appl. 14(3), 703-714 (2008)
  \item Bracken, C., Byrne, E., Markin, N., McGuire, G.: Determining the nonlinearity of a new family of APN functions. In: Boztas, S., Lu, H.-F. (eds.) Proc. AAECC-17 Conference. LNCS, vol. 4851, pp. 72-79 (2007)
  \item Budaghyan, L., Carlet, C.: Classes of quadratic APN trinomials, hexanomials and related structures. IEEE Trans. Inf. Theory 54(5), 2354-2357 (2008)
  \item Budaghyan, L., Carlet, C., Felke, P., Leander, G.: An infinite class of quadratic APN functions which are not equivalent to power mappings. In: Proceedings of ISIT 2006, Seattle, USA (2006)
  \item Budaghyan, L., Carlet, C., Pott, A.: New constructions of almost bent and almost perfect nonlinear functions. IEEE Trans. Inf. Theory 52(3), 1141-1152 (2006)
  \item Budaghyan, L., Carlet, C., Leander, G.: Another class of quadratic APN binomials over $F_{2^{n}}$ : the case $n$ divisible by 4. In: Proceedings of WCC 07, pp. 49-58. Versailles, France (2007)
  \item Budaghyan, L., Carlet, C., Leander, G.: Two classes of quadratic APN binomials inequivalent to power functions. IEEE Trans. Inf. Theory 54(9), 4218-4229 (2008)
  \item Budaghyan, L., Carlet, C., Leander, G.: Constructing new APN functions from known ones. Finite Fields Their Appl. 15(2), 150-159 (2009)
  \item Carlet, C., Charpin, P., Zinoviev, V.: Codes, bent functions and permutations suitable for DESlike cryptosystems. Designs Codes Cryptogr. 15(2), 125-156 (1998)
  \item Dillon, J.: Slides from talk given at Polynomials over Finite Fields and Applications. Held at Banff International Research Station (2006)
  \item Edel, Y., Kyureghyan, G., Pott, A.: A new APN function which is not equivalent to a power mapping. IEEE Trans. Inf. Theory 52(2), 744-747 (2006)
  \item Feulner, T.: APN functions, available online at \href{http://www.algorithm.uni-bayreuth.de/en/research/}{http://www.algorithm.uni-bayreuth.de/en/research/} Coding\_Theory/APN\_Functions/index.html
  \item Nyberg, K.: Differentially uniform mappings for cryptography. Advances in CryptologyEUROCRYPT 93. Lecture Notes in Computer Science, pp. 55-64. Springer-Verlag (1994)
\end{enumerate}


\end{document}